\begin{textsample}{NEG Dim 1 – ba – Score 6.00 – ba/ba\_inf\_e\_ef\_22.txt}  \label{ex:f1_neg_190}
5.6.
Educação Financeira e para o Consumo A sociedade contemporânea vive um momento de crise, em que se faz necessária a mudança do paradigma antropocêntrico.
Os padrões de consumo impostos pela “ sociedade ”, por meio do sistema econômico predominante, devem ser revistos, sob pena de inviabilizar a continuidade da vida no planeta.
A educação possui papel fundamental na formulação de uma nova mentalidade, e a Educação Financeira e para o Consumo é elemento-chave na formação de uma consciência em relação à responsabilidade social na busca da qualidade de vida das pessoas e do planeta.
Em uma sociedade em que é mais importante o TER do que o SER, abrem -se as portas para a discussão sobre o consumo consciente e sobre o que, como e por que consumimos.
Neste contexto, o Tema Integrador Educação Financeira e para o Consumo visa a construção e o desenvolvimento de comportamentos financeiros consistentes, autônomos e saudáveis, para que os estudantes possam, como protagonistas de suas histórias, planejar e executar os seus projetos de vida.
Ferreira ( 2017 ), em seu artigo intitulado “ A importância da educação financeira pessoal para a qualidade de vida ”, apresenta argumentos e relaciona os índices de qualidade de vida com os conhecimentos e práticas da educação financeira pessoal, destacando que não há intenção de : > “ [... ] expor que qualidade de vida é parar de gastar ou poupar apenas para item específico, e sim mostrar que gastando de forma consciente e inteligente o indivíduo tem mais possibilidade de conquistar o que para ele é importante, assim como proporcionar uma vida mais tranquila e estável sem um endividamento constante que acaba por tirar a tranquilidade do indivíduo. ” As unidades escolares devem promover a inserção de conteúdos que estimulem a capacidade de escolha consciente e responsável nas discussões em sala de aula, apontando para a formação de indivíduos que possam gerir/mediar os recursos, transcendendo a questão restrita ao dinheiro, ou seja, não versado na aquisição de bens associados, tão somente, ao lucro imediato, mas para a constituição de cidadãos que reconheçam o caráter finito dos recursos e, portanto, capazes de agregar bens sem desconsiderar o \textbf{desperdício} e o descarte irresponsável destes no ambiente e, principalmente, o consumismo desenfreado.

% matched lemmas: desperdício
\end{textsample}
